\documentclass{article}
\usepackage[margin=2cm]{geometry}
\usepackage{amsmath}
\usepackage{graphicx}
\usepackage{booktabs}
\usepackage[utf8]{inputenc}
\usepackage{ragged2e}
\setlength{\emergencystretch}{3em}
\begin{document}
\sloppy

\noindent\hfill\includegraphics[height=1.8em]{assets/figure_header_logo.png}\par\noindent\hrule\vspace{4pt}

expectations gives


\begin{equation}
\begin{array} { r } { F _ { Y _ { 1 } ( 1, 0 ) | D = 0, M ( 0 ) = 0 } ( y ) = \frac { p _ { n } } { p _ { n } + p _ { c } } F _ { Y _ { 1 } ( 1, 0 ) | D = 0, M ( 0 ) = 0, \tau = n } ( y ), } \\ { \qquad + \frac { p _ { c } } { p _ { n } + p _ { c } } F _ { Y _ { 1 } ( 1, 0 ) | D = 0, M ( 0 ) = 0, \tau = c } ( y}\end{array}
\end{equation}


Using () and rearranging the equation gives,


\begin{equation}
F _ { Y _ { 1 } ( 1, 0 ) | r = c } ( y ) = \frac { p _ { 0 | 0 } } { p _ { 0 | 0 } - p _ { 0 | 1 } } F _ { Q t s ( Y _ { 0 } ) | D = 0, M = 0 } ( y ) - \frac { p _ { 0 | 1 } } { p _ { 0 | 0 } - p _ { 0 | 1 } } F _ { Y _ { 1 | D = 1, M = 0 } } ( y ). ~
\end{equation}


In analogy to Q, the outcome distribution under $D = 0$and$i i = 0$equals:


\begin{equation}
F _ { Y _ { 1 } | D = 0, M = 0 } ( y ) = \frac { p _ { n } } { p _ { n } + p _ { c } } F _ { Y _ { 1 } ( 0, 0 ) | r = n } ( y ) + \frac { p _ { c } } { p _ { n } + p _ { c } } F _ { Y _ { 1 } ( 0, 0 ) | \tau = c } ( y ).
\end{equation}


Using () and rearranging the equation gives


\begin{equation}
E _ { Y _ { 1 } ( \mathbf { 0 }, \mathbf { 0 } ) | r = c } ( y ) = \frac { p _ { 0 | 0 } } { p _ { 0 | 0 } - p _ { 0 | 1 } } F _ { Y _ { 1 } | D = 0, M = 0 } ( y ) - \frac { p _ { 0 | 1 } } { p _ { 0 | 0 } - p _ { 0 | 1 } } F _ { Q _ { 0 | 1 } | D = 1, M = 0 } ( y ). ~
\end{equation}


\subsection*{1  Proof of Theorem 4.}


\subsection*{1.1 Average direct effect on the always-takers}


In the following, we show that$\theta . { 1 } ^ { a } = E [ Y . { 1 } ( 1, 1 ) \! - \! Y . { 1 } ( 0, 1 ) | \tau = a ] = E [ Q . { 1 1 } ( Y . { 0 } ) \! - \! Y . { 1 } | D =$
$0, \mathtt { i } \mathbf { j } = 1 ]$. From , we obtain the first ingredient$\left ( \right )$ $E [ Y . { 1 } ( 0, 1 ) | a ] = E [ Y . { 1 } | D = 0, ~ M =$
1].Furthermore, from $( )$ we  have $E [ Q . { 1 1 } ( Y . { 0 } ) | D ~ = ~ 0, M ~ = ~ 1 ] ~ = ~ E [ Y . { 1 } ( 1, 1 ) | D ~ =$
$0. i I ( 0 ) = 1 ]$ Under Assumption 7 and 8,.

\end{document}
